We represent a manuscript as a directed, labeled document graph $G=(V,E,L,M)$: each vertex $v\in V$ is an atomic writing unit with metadata $m\in M$ (id, title, target length, optionality, evidence\_scope, generation constraints), and each edge $e=(u\!\rightarrow\!v)\in E$ is labeled by $\ell\in L$ (e.g., supports, requires) indicating that $u$ must be available when generating $v$.  
We require $G$ to be acyclic; cycles trigger cycle‑breaking suggestions and the scheduler uses a reverse‑topological (leaf‑first) traversal so evidence‑bearing nodes are produced before higher‑level synthesis.  
Scheduling combines hard precedence with soft heuristics (evidence‑availability, node criticality/out‑degree, synthesis penalty, estimated write cost); at each step the ready set $R$ is scored by $P(v)$ (a weighted sum) and $v^*=\arg\max_{v\in R}P(v)$ is generated (with evidence injection and JSON constraints) until completion.  
Optional nodes use a skip\_threshold; cross‑cutting sections are handled either by explicit incoming edges or by a final synthesis pass seeded with structured pointers; failures trigger local regeneration, retrieval expansion, or marking unresolved with diagnostic tags, and a provenance trace (node id, timestamp, evidence ids, status) is recorded in the JSON outputs.  
These choices are heuristic and configurable; implementation details are in Section 4 and failure analyses in Sections 3.4 and 6.2.